\documentclass[11pt]{article}
\usepackage{amsmath,amssymb,amsthm,mathtools,hyperref}

\begin{document}
\section*{Linking microevolution to macroevolution}
\label{prob:160}
\label{rank:33}
\par\smallskip

\noindent\textbf{Domain.} Speciation; macroevolutionary coarse graining. \quad \textbf{Retained status/source context.} Open cross-scale synthesis problem. Comparative speciation research motivates it, but does not assert that a unique diversification process can be recovered from population dynamics.
\paragraph{Definitions and assumptions.} Let populations occupy a finite habitat network. A forward stochastic model specifies birth, death, migration, mutation and mate acceptance as functions of genotype and environment. Define normalized gene exchange $R_{ij}(t)$ as expected realized gene exchange between populations $i,j$, normalized to a freely interbreeding control. A species split occurs only if $R_{ij}<\epsilon$ persists for $H$ generations; extinction means loss of all individuals in a species. Once an initial species partition and a complete partition and lineage update rule satisfying these event conditions have been declared, they define a species history $\mathcal T$ and species count $S_t$. The approximation compares a fixed horizon and event resolution rather than treating all geological times equally. Sensitivity to the isolation threshold and persistence duration is reported, because these choices can turn transient divergence into an apparent speciation event.
\paragraph{Problem statement.} Derive an approximation by speciation and extinction hazards $\lambda,\mu$ with an explicit error criterion:
\[
 \sup_{0\le t\le T}
 \left|\Pr_{\mathrm{micro}}(S_t=k)-\Pr_{\lambda,\mu}(S_t=k)\right|
 \le\delta\quad\text{for every }k\in\{0,\ldots,K\}.
\]
Identify timescale-separation and sampling conditions under which the hazards depend only on declared ecological covariates, rather than unrecorded population history.
\paragraph{Instance specification.} The binding catalogue conventions apply to both probability laws in the displayed comparison. A concrete instance declares the microscopic model family and initial laws, environmental and validation domain, time convention and units, horizon $T$, nonnegative integer cutoff $K$, positive tolerances and normalization scales, isolation threshold $\epsilon$, and persistence duration $H$ in generations, including any isolation history needed at the initial time. For $R_{ij}(t)$, declare the exchanged quantity and its units, directional aggregation, exchange window, expectation law and conditioning information, and freely interbreeding control, with a positive-denominator or explicit undefined/zero convention. Normalization alone does not imply $R_{ij}(t)\le1$. Declare the complete partition and lineage update rule, including simultaneous events, reconnection and persistence-timer updates. Define the comparison law by its state, initial law, transition mechanism, rate units and allowed hazard dependencies; specify whether hazards are per-lineage or total rates and how the comparison handles any mergers or non-unit count jumps admitted by the microscopic species rule. The stochastic laws must satisfy the applicable well-posedness and non-explosion requirements below. Declare the ecological covariates, their dynamics and the conditioning information for the hazards, and whether history-independent closure is exact or approximate; any approximation claim must state how its error is assessed, including the treatment of persistence timers and other hidden state. State the order of quantifiers over microscopic laws, initial laws, environments and candidate hazards, any timescale regime, and the allowed dependence of hazards and error bounds on the declared inputs. Specify the role of sampling and distinguish complete species counts from sampled or reconstructed counts wherever observations are used. These declarations are instance inputs or conditions to be identified, not conventions selected by this scheme.
\paragraph{Additional binding assumptions and resolution criterion.} The quantity $R_{ij}$ is normalized gene exchange, so smaller values indicate greater isolation. Pairwise thresholding is not transitive and therefore does not uniquely define species: specify a partition and split/merge update rule, including reconnection during the persistence interval. Bounds for $k\leq K$ do not control the unobserved count tail; add a tail bound or use total variation on the full count law when claiming a distributional approximation. The displayed criterion compares one-time count marginals uniformly over the declared time domain; even with full count-tail control, it does not compare joint count trajectories or species-history laws.

\paragraph{Alternative formulations.} (i) Under rare speciation and rapid within-population equilibration, test a Markov birth--death approximation. (ii) When isolation is reversible or hybridization occurs, use a history-dependent network process rather than a bifurcating species tree. These models need not have equivalent diversification rates.
\paragraph{Obstructions and partial progress.} Population and comparative studies explain components of reproductive isolation, and many macroevolutionary birth--death models can be fitted directly. Their rates need not correspond uniquely to microevolutionary mechanisms, especially when extinct or unsampled lineages are absent. The unresolved task is justified coarse graining and predictive transfer of rates, including sensitivity to the operational species criterion.
\paragraph{Sources (supplied provenance; not all rechecked).}
\begin{itemize}
\item Matute and Cooper, Comparative studies on speciation: 30 years since Coyne and Orr (2021). \url{https://pmc.ncbi.nlm.nih.gov/articles/PMC8247902/}. Supplied review; indexed metadata/abstract accessed, direct full-text request challenged.
\item Matute and Cooper, Comparative studies on speciation: 30 years since Coyne and Orr (2021). \url{https://pubmed.ncbi.nlm.nih.gov/33491225/}. Review metadata; comparative research agenda rather than a macroevolutionary derivation.
\end{itemize}


\section*{Binding catalogue conventions}
This edition revises both supplied catalogues and preserves original problem numbers 1--330. It adds one globally unique label \verb|\label{rank:x}| per problem, where smaller $x$ means greater provisional difficulty. A rank label is an anchor at that problem's section: \verb|\ref{rank:x}| returns the section number, not the literal difficulty number. The visible \texttt{rank:x} text and the index give the difficulty number itself.

The entries are \emph{parameterized mathematical research schemes}, including empirical prediction, causal identification, historical reconstruction, theory and engineering. They are not 330 certified unsolved mathematical conjectures. A rigorous instance fixes the objects and restrictions below; changing those objects changes the question. Established identities, finite calculations, non-identification constructions and solved special cases are not being posed anew as open. Some residual tasks concern informative bounds, biological realization or external validation rather than a new theorem.

Every main definition and problem statement received a mathematical editorial review. Each entry now has an additional binding-conditions paragraph; direct corrections are separately recorded in the accompanying audit. The alternative formulations inherit the applicable conventions, but are different targets unless equivalence is proved. Where an old broad phrase conflicts with an explicit restriction below it, the explicit restriction controls.

\textbf{Literature boundary.} Thirteen selective primary-literature checks cover eighteen entries in this revision; their access levels and scope are recorded explicitly. The supplied bibliography and its earlier source-access descriptions are retained as provenance, not represented as newly verified. Neither this review nor a source's description proves the absence of later solutions. The date of this edition is 13 October 2026; there is no claim of exhaustive current-status certification for all 330 entries.

\subsection*{Binding conventions for a concrete instance}
\paragraph{Objects, quantifiers and non-vacuity.}
An instance declares measurable state, input, output and parameter spaces; a fixed candidate family; admissible initial laws, interventions and environments; a finite horizon or specified asymptotic regime; observation and sampling laws; and performance metrics, positive normalization scales and positive tolerances. These are chosen before fitting or testing, not adjusted after seeing failures. Shared parameters must genuinely be shared across the declared validation domain. Finite dimension alone does not restrict arbitrary real-number encodings: specify a finite library, bounded description length, or regularity/complexity bounds. Exclude constant or no-signal solutions through a fixed informative-output or baseline-improvement requirement where appropriate. A universally quantified result cannot be established by choosing one trivial favorable instance.

\paragraph{Probability laws and data.}
Every expectation or probability refers to the specified stochastic law, including initial uncertainty, process noise, observation noise and sampling as appropriate. A deterministic equation is not a probability model without a declared random input or initial law. ``Observed'' targets in population-level formulas mean population observables $g_j^*$, not noisy estimates asserted to be exact. Let $\widetilde g_j$ be independent validation estimates and suppose
\[
\mathbb{P}\{\lvert\widetilde g_j-g_j^*\rvert\le r_j\text{ for every }j\}\ge1-\alpha.
\]
Then an uncertainty-aware sufficient certification rule for positive target tolerances $\varepsilon_j$ is
\[
\max_j\frac{|\widehat g_j-\widetilde g_j|+r_j}{\varepsilon_j}\le1.
\]
On that simultaneous confidence event it implies every population error is at most its tolerance, by the triangle inequality. Vector/norm versions use corresponding valid simultaneous bounds. Training uncertainty must also be included when the claim averages over repeated training data. Evaluation on finitely many conditions does not establish a continuous-domain supremum without covering and regularity bounds.

\paragraph{Dynamics and well-posedness.}
For ODEs and SDEs give coefficient domains, initial laws and existence/uniqueness conditions in the intended solution sense, or keep the entire admitted solution set. For stochastic jump models specify nonnegative off-diagonal rates, diagonal normalization, rate orientation and non-explosion. For PDEs specify spatial domains, boundary/initial conditions, conservation or source terms, and a weak or classical solution class. Moving-domain concentration equations include dilution/transport terms. A stationary law is used only for a time-homogeneous augmented model under existence assumptions; a stationary approximation to a driven system must be justified. Absorbing extinction states are retained. Survival-conditioned, quasi-stationary and unconditional laws are distinct and must not be substituted for one another.

\paragraph{Identification and causal meaning.}
Two parameter values are observationally equivalent when they induce the same full distribution under all allowed observation/intervention designs, not merely the same likelihood value for one sample. A target is identified only if constant on the equivalence class. Gauge symmetries, label switching and reference conventions are quotiented or fixed. Finite-data compatible sets use a declared confidence, likelihood or discrepancy rule; an empty set signals model rejection. Bayesian credible sets are not automatically confidence sets. A prior or regularizer does not create data identification. Bounds intended to resolve a question must be informative at a predeclared width or decision threshold.

A causal contrast uses well-defined intervention versions and a stated causal model. Random assignment, or explicit consistency, overlap and exchangeability assumptions, is required for causal interpretation from observations. Interference and time-varying confounding are represented when present. Fixing a mediator by intervention differs from conditioning on its observed value. Marginal potential-outcome laws do not determine individual cross-world effects or an arbitrary joint coupling. Ethically admissible designs, rather than unconstrained interventions, define the experimental menu.

\paragraph{Metrics, units, optimization and limits.}
Specify the sampling or reference measure for densities and integrals, the log base, and dimensionless arguments of logarithms. Total variation is $\sup_A|P(A)-Q(A)|$; Wasserstein distances require the declared ground metric and appropriate moment conditions; KL divergence uses its stated absolute-continuity convention. All normed readouts are scaled consistently, and error probabilities use $0<\alpha<1$. A ratio has a positive denominator or an explicit undefined/zero convention. Thermodynamic entropy production is nonnegative in expectation under the relevant hypotheses; individual stochastic realizations need not be. An infimum is not an attained minimum without compactness/coercivity and continuity or another existence argument. Fix the order of large-population, long-time and vanishing-mutation limits. Stability claims declare perturbation space and symmetries.

\subsection*{How the provisional difficulty order was assigned}
Four editorial components are scored from 1 (lower burden) to 5 (higher burden): (a) formal or inferential difficulty, (b) access to evidence and identification burden, (c) mechanistic/computational integration, and (d) independent validation and transfer. Their unweighted total orders the entries. For selected entries in the highest bands, an explicit editorial priority breaks equal totals; all other ties are broken by original problem number. Equal totals are \emph{not} evidence of reliably different true difficulties. Scores are judgments, not measurements, completion-time estimates or consensus rankings.

The object ranked is an informative, nondegenerate resolution of the \emph{scoped} scientific target, with the stated budget and validation domain. A sharp useful impossibility result can qualify; an already known identity, vacuous bound, arbitrary favorable model or lookup table cannot. Rank 330 is merely the least difficult relative assessment in this catalogue, not a claim of an easy biological solution. Broad historical titles do not automatically outrank precisely formulated mechanistic tasks.



\end{document}
